\section*{程甲本}
話說尤三姐自戕之後，尤老娘以及尤二姐、賈珍、尤氏並賈蓉、賈璉等俱不勝悲慟傷感，忙着買棺盛殮，送往城外埋葬。柳湘蓮見尤三姐身亡，迷性不悟，尚有癡情眷戀，卻被道人數句偈言打破迷關，竟自削髮出家，隨一瘋道人飄然而去，不知何往。

薛姨媽聞知湘蓮已說定了尤三姐，正打算替他買房置器，擇日迎娶過門，以報他救命之恩。忽有家中小廝吿知尤三姐自戕與柳湘蓮出家之事，心甚嘆息。時值寶釵從園中過來，聽了這些話，並不在意，乃勸道：「俗語說的好，『天有不測風雲，人有旦夕禍福』。這也是他們前生命定。前日媽媽為他救了哥哥，商量著替他料理，如今已經死的死了，走的走了，依我說，也只好由他罷了。媽媽也不必為他們傷感了。倒是自從哥哥打江南回來了一二十日，販了來的貨物，想來也該發完了，媽媽和哥哥商議商議，酬謝酬謝那同去的那張德輝才是。伙計們辛辛苦苦的，回來幾個月了，也該請一請，別叫人家看著無禮似的。」

母女正說話間，見薛蟠自外而入，眼中尚有淚痕，一進門來，便向他母親拍手說道：「媽媽可知道柳二哥尤三姐的事麼？」薛姨媽說：「我才聽見說，正在這裡合你妹妹說這件公案呢。」薛蟠道：「媽媽可聽見說湘蓮跟著一個道士出了家了麼？」薛姨媽道：「這越發奇了。怎麼柳相公那樣一個年輕的聰明人，一時胡塗了，就跟著道士去了呢？我想你們好了一場，他又無父母兄弟，隻身一人在此，你該各處找找他才是。靠那道士，能往那裡遠去？左不過是在這方近左右的廟裡寺裡罷了。」薛蟠說：「何嘗不是呢？我一聽見這個信兒，就連忙帶了小廝們在各處尋找，連一個影兒也沒有。又去問人，都說沒看見。」

薛姨媽說：「你既找尋過，沒有，也算把你做朋友的心盡了。再者，你妹妹才說你也回家半個多月了，想貨物也該發完了，也該擺桌酒，給張德輝和夥計們，道道乏才是。」薛蟠聽說，便道：「媽媽說的很是。倒是妹妹想的周到。因這些日子，為各處發貨，又為柳二哥的事忙了這幾日，把正經事都誤了。要不然，定了明兒後兒，下帖兒請罷。」薛姨媽道：「由你辦去罷。」

話猶未了，外面小廝在門外回說：「張管總着人送了兩個箱子來。」薛蟠聽了，便命小廝央門外幾個夥計搬進了兩個夾板夾的大棕箱。薛蟠一見說：「特給媽和妹妹帶來的東西，不是夥計送家裡來，我都忘了。」

薛姨媽同寶釵問：「是什麼好東西，這樣捆着夾着的？」便命人挑了繩子，去了夾板，開了鎖看時，卻是些綢緞、綾錦、洋貨等家常應用之物。獨有寶釵他的那個箱子裏，除了筆、墨、硯、各色箋紙、香袋、香珠、扇子、扇墜、花粉、胭脂、頭油等物外，還有虎丘帶來的自行人、酒令兒、水銀灌的打筋斗的小小子，沙子燈，一出一出的泥人兒的戲，用青紗罩的匣子裝着，又有在虎丘山上作的薛蟠的小像，泥捏成的與薛蟠毫無相差，以及許多碎小頑意兒的東西。寶釵一見，拿着薛蟠的小像細細看了，又看看他哥哥捂着嘴微笑，再和母兄說了一回閒話。便吩咐鶯兒：「你帶幾個老婆子，將我的這個箱子，拿到園子裏去，我好就近從那邊送人。」說着，便起身辭了母兄往園子裏去了。這裏薛姨媽將自己這個箱子裏的東西取出，一份一份的打點清楚，着鶯兒送往賈母並王夫人等處。

寶釵隨着箱子到了自己房中，將東西逐件過了目，除將自己留用之外，遂一一配妥當：也有送筆、墨、紙、硯的，也有送香袋、扇子、香墜的，也有送脂粉、頭油的，有單送頑意兒的。一一打點完畢，使鶯兒同一個老婆子，送往各處。

寶釵送東西的ㄚ頭回來，說：「也有道謝的，也有賞錢的，獨有給巧姐兒的那一份，仍舊拿回來了。」寶釵一見，不知何意，便問：「為什麼這一份沒送去，還是送了去沒收呢？」鶯兒說：「我方才給環哥兒送東西的時候，見璉二奶奶往老太太房裡去了。」寶釵說：「二奶奶不在家，你只管交給丫頭們收下，等二奶奶回來，自有他們告訴就是了。」鶯兒聽了，又與老婆子出了園子，到了鳳姐這邊，送了東西，回來見寶釵。

寶釵問道：「你見了璉二奶奶沒有？」鶯兒說：「我沒見。」寶釵說：「二奶奶還沒有回來？」鶯兒說：「回來是回來了。因豐兒對我說：『二奶奶自老太太屋裡回來，一臉怒氣，叫了平兒去，唧唧咕咕的說話，也不叫人聽見。你不必見，等我替你回一聲兒就是了。』因此豐兒拿進去，回了二奶奶。我們就回來了。」寶釵聽了，自己納悶，想不出鳳姐是為什麼生氣。

衆人不過收了東西，皆說些見面再謝等語而已。惟有林黛玉見是江南家鄉之物，便對着揮淚自嘆。紫鵑深知黛玉心腸，在一旁勸道：「寶姑娘送來這些東西，姑娘看着該喜歡才是。」

話猶未畢，只見寶玉已進來。寶玉見黛玉淚痕滿面，便問：「妹妹，又是為的什麼？」黛玉不答。旁邊紫鵑將嘴向牀後桌上一努，寶玉會意，便往牀上一看，見堆着許多東西，就知道是寶釵送來的。寶玉深知黛玉是因見了江南來的故鄉之物，勾起傷感落淚。便道：「妹妹，你放心！等我明年往江南去，與你帶兩船來。」黛玉聽了這話，說道：「你那裏知道我的緣故。」說着眼淚又流了下來。寶玉忙走到牀前，挨著黛玉坐下，將那些東西一件一件拿起來，擺弄著細瞧，故意問:「這是什麼，叫什麼名字?那是什麼做的，這樣齊整?這是什麼，要它做什麼使用?妹妹，你瞧，這一件可以擺在書閣兒上作陳設，那件放在條案上當古董兒倒好呢!」一味的將些沒要緊的話來支吾，搭訕。黛玉見寶玉可笑的樣子，稍將煩惱丟開。寶玉便說道:「寶姐姐送東西來給咱們，我想著，咱們也該到她那里道個謝去才是，不知妹妹可去不去?」黛玉道：「自家姐妹，這倒不必。只是到他那邊，薛大哥回來了，必然告訴他些南邊的新聞故事兒，我去聽聽，只當回了家鄉一趟的。」說著，眼圈兒又紅了。寶玉便站著等他。黛玉只得和他出來，往寶釵那裡去了。

二人到寶釵處，道了謝，寶玉又口口稱贊泥人兒等物有趣。寶釵笑道：「原不是什麼好東西，不過是遠路帶來的土物兒，大家看著新鮮些就是了。」黛玉道：「這些東西我們小時候倒不理會，如今看見，真是新鮮物兒了。」寶釵因笑道：「妹妹知道，這就是俗語說的『物離鄉貴』，其實可算什麼呢。」寶玉聽了這話正触著黛玉方纔的心事，連忙拿話岔開：「明年大哥哥還去江南嗎？——」話沒說完，黛玉早接口道：「——姐姐，你瞧，寶哥哥不是給姐姐來道謝，竟又要定下明年的東西來了。」說的寶釵寶玉都笑了。

三個人又閒話了一回，因提起黛玉的病來，寶釵勸了一回，因說道：「妹妹若覺著身上不爽快，倒要自己勉強扎掙著出來，各處走走逛逛，散散心，比在屋裡悶坐著到底好些。我那兩日，不是覺著發懶，渾身發熱，只是要歪著？也因為時氣不好，怕病，因此尋些事情，自己混著。這兩日才覺得好些了。」黛玉道：「姐姐說的何嘗不是？我也是這麼想著呢。」大家又坐了一會子方散。寶玉仍把黛玉送至瀟湘館門首，才各自回去了。

且說那趙姨娘因見寶釵送環哥兒物件，心中甚喜，滿嘴誇獎：「人人都說寶姑娘會行事，很大方，今日看來，果然不錯。他哥哥能帶了多少東西來，他挨家送到，並不遺漏一處，也不露出誰薄誰厚，連我們他都想到了，若是林姑娘，即或有人帶了東西來，那裏輪得到我們娘兒倆身上呢！可見人會行事，真真露着各別另樣的好。」趙姨娘因環哥兒得了東西，深為得意，不住的托在掌上擺弄瞧看一會。想寶釵乃係王夫人之表侄女，特要在王夫人跟前賣好兒。自己蝎蝎螫螫的拿着那東西，走至王夫人房中，站在一旁說道：「這是寶姑娘才給環哥的，他年輕輕的人想得周到，我還給了送東西的小ㄚ頭二百錢。聽說姨太太也給太太送來了，不知是什麼東西？你們瞧瞧這一個門裏頭，就是兩份兒，能有多少呢？怪不得老太太同太太都誇他疼他，果然招人疼。」說着，將手裏的東西遞過去與王夫人瞧，誰知王夫人頭也沒擡，手也沒伸，只口內說了聲「好，給環哥兒頑去罷」，並無正眼看一看。趙姨娘因招了一鼻子灰，滿肚氣惱，無精打彩的回房，將東西丟在一邊，也無人問他，他卻自己咕嘟着嘴，一邊子坐着。

且說薛蟠聽了母親之言，次日請了張德輝與四位夥計，俱已到齊，不免說些販賣賬目發貨之事。不一時，上席讓坐，薛蟠挨次斟了酒，薛姨媽又使人出來致謝，大家喝著酒說閒話兒。內中一個道：「今兒這席上短了柳二爺。」薛蟠聞言，把眉一皺，嘆口氣道：「什麼是柳二爺，如今不知那裡作『柳道爺』去了。」眾人都詫異道：「這是怎麼說？」薛蟠便把湘蓮前後事體說了一遍。眾人聽了，越發駭異，因說道：「怪不的。前兒我們在店裡，髣髣髴髴也聽見人吵嚷，說：『有一個道士，三言兩語，把一個人度了去了。』又說「『一陣風颳了去了。』只不知是誰。我們正發貨，那裡有閒工夫打聽這個事去？到如今還是似信不信的，誰知就是柳二爺呢？張德輝道：「柳二爺那樣個伶俐人，未必是真跟了道士去罷。他原會些武藝，又有力量，或看破那道士的妖術邪法，特意跟他去，在背地擺佈他，也未可知。」薛蟠道：「果然如此，倒也罷了。」眾夥計隨便喝了幾杯酒，吃了飯，大家散了。

話說寶玉回來，想着黛玉的孤苦，不免替他傷感起來。襲人見寶玉從外面進來坐在那發呆，便問：「就回來了？是不是同林姑娘一塊去了寶姑娘那兒？」寶玉說：「我會林姑娘同去的——送林姑娘的東西比送我們的多一兩倍呢。」說着話兒，便叫取了枕來，要在牀上歪着。襲人說：「璉二奶奶自從病了一場之後，我早就想着要到他那裏去看看，你同晴雯麝月呆着，我去看看就來。」寶玉說：「你只管去罷。」言畢，襲人遂換了兩件新鮮衣服。囑咐了晴雯、麝月幾句，便出了怡紅院。

至沁芳橋上立住，往四下里觀看那園中景致。那時正是夏末秋初，園內蟬鬧蟲鳴；只是花也開敗了，芙蓉池中荷葉新殘相間，也將殘上來了。倒是近着池邊，都發了紅鋪鋪的咕嘟子，襯着碧綠的葉兒，着實可愛。於是一壁里瞧着，一壁里下了橋。走了不遠，迎見李紈房裏的丫頭素雲捧着個洋漆盒兒走來。襲人便問：「往那裏去送東西？」素雲說：「這是我們奶奶給三姑娘送去的菱角兒、雞頭米。」襲人說：「這個東西，是咱們園子裏河內采的，還是外頭買來的呢？」素雲說：「是我們那邊劉媽媽的女兒從鄉下帶來孝敬我們奶奶的。因三姑娘在我們那裏坐，奶奶叫人剝了讓他吃。他說：『才吃了熱茶了，一會子再吃罷。』所以命我給三姑娘送過去。」言畢，各自散了。

襲人走著，沿堤看頑了一回。猛抬頭看見那邊葡萄架底下有人拿著撣子在那里撣什麼呢，走到跟前，卻是老祝媽。那老婆子見了襲人，便笑嘻嘻的迎上來，說道：「姑娘怎麼今日得工夫出來逛逛？」襲人道：「可不是。我要到璉二奶奶家去。你在這里做什麼呢？」那婆子道：「我在這里赶蜜蜂兒。今年三伏里雨水少，這果子樹上都有虫子，把果子吃的疤瘌流星的掉了好些下來。姑娘還不知道呢，這馬蜂最可惡的，一嘟嚕上只咬破三兩個兒，那破的水滴到好的上頭，連這一嘟嚕都是要爛的。姑娘你瞧，咱們說話的空兒沒赶，就落上許多了。」襲人道：「你就是不住手的赶，也赶不了許多。你倒是告訴買辦，叫他多多做些小冷布口袋兒，一嘟嚕套上一個，又透風，又不遭塌。」婆子笑道：「倒是姑娘說的是。我今年才管上，那里知道這個巧法兒呢。」

襲人說：「如今這園子裡這些果品有好些種，到是那樣先熟的快些?」老祝婆子說：「如今才入七月的門，果子都是才紅上來，要是好 吃，想來還得月盡頭兒才熟透了呢。姑娘不信，我摘一個給姑娘嚐嚐。」襲人正色說道:「這那裡使得?不但沒熟吃不得，就是熟了，一則沒有 供鮮，二則主子們尚然沒吃，我如何先吃得呢?」老婆子忙笑道：「姑娘說得有理。我因為姑娘問我，我白這樣說。」襲人說：“我方才告訴你要口袋的話，你就回一回二奶奶，叫管事的作去罷。」言畢，遂一直的出了園子的門，就到鳳姐這裡來了。

 一到院里，只聽鳳姐說道：「天理良心，我在這屋里熬的越發成了賊了。」襲人聽見這話，知道有原故了，又不好回來，又不好進去，遂把腳步放重些，隔著窗子問道：「平姐姐在家里呢么？」平兒忙答應著迎出來。襲人便問：「二奶奶也在家里呢么，身上可大安了？」說著，已走進來。鳳姐裝著在床上歪著呢，見襲人進來，也笑著站起來，說：「好些了，叫你惦著。怎麼這幾日不過我們這邊坐坐？」襲人道：「奶奶身上欠安，本該天天過來請安才是。但只怕奶奶身上不爽快，倒要靜靜兒的歇歇兒，我們來了，倒吵的奶奶煩。」鳳姐笑道：「常聽見平兒說你背地里還惦著我，常常問我。這就是你盡心了。」一面說著，叫平兒挪了張杌子放在床旁邊，讓襲人坐下。丰兒端進茶來，襲人欠身道：「妹妹坐著罷。」一面說閒話兒。只見一個小丫頭子在外間屋里悄悄的和平兒說：「旺兒來了。在二門上伺候著呢。」襲人知他們有事，又說了兩句話，便起身要走。鳳姐道：「閒來坐坐，說說話兒，我倒開心。」因命平兒：「送送你妹妹。」平兒答應著送出來。只見兩三個小丫頭子，都在那里屏聲息氣齊齊的伺候著。襲人不知何事，便自去了。

卻說平兒送出襲人，進來回道：「旺兒才來了，因襲人在這裡我叫他先到外頭等等兒，這會子還是立刻叫他呢，還是等著？請奶奶的示下。」鳳姐道：「叫他來。」平兒忙叫小丫頭去傳旺兒進來。這裡鳳姐又問平兒：「你到底是怎麼聽見說的？」平兒道：「就是頭裡那小丫頭子的話。他說他在二門裡頭聽見外頭兩個小廝說：『這個新二奶奶比咱們舊二奶奶還俊呢，脾氣兒也好。』不知是旺兒還是誰，吆喝了兩個一頓，說：『什麼新奶奶舊奶奶的，還不快悄悄兒的呢，叫裡頭知道了，把你的舌頭還割了呢。』」平兒正說著，只見一個小丫頭進來回說：「旺兒在外頭伺候著呢。」鳳姐聽了，冷笑了一聲說：「叫他進來。」那小丫頭出來說：「奶奶叫呢。」旺兒連忙答應著進來。旺兒請了安，在外間門口垂手侍立。鳳姐兒道：「你過來，我問你話。」旺兒才走到裡間門旁站著。鳳姐兒道：「你二爺在外頭弄了人，你知道不知道？」旺兒又打著千兒回道：「奴才天天在二門上聽差事，如何能知道二爺外頭的事呢。」鳳姐冷笑道：「你自然不知道。你要知道，你怎麼攔人呢。」旺兒見這話，知道剛纔的話已經走了風了，料著瞞不過，便又跪回道：「奴才實在不知。就是頭裡興兒和喜兒兩個人在那裡混說，奴才吆喝了他們兩句。內中深情底裡奴才不知道，不敢妄回。求奶奶問興兒，他是長跟二爺出門的。」鳳姐聽了，下死勁啐了一口，罵道：「你們這一起沒良心的混帳忘八崽子！都是一條藤兒，打量我不知道呢。先去給我把興兒那個忘八崽子叫了來，你也不許走。問明白了他，回來再問你。好，好，好，這才是我使出來的好人呢！」那旺兒只得連聲答應幾個是，磕了個頭爬起來出去，去叫興兒。

卻說興兒正在帳房兒里和小廝們玩呢，聽見說二奶奶叫，先唬了一跳，卻也想不到是這件事發作了，連忙跟著旺兒進來。旺兒先進去，回說：「興兒來了。」鳳姐兒厲聲道：「叫他！」那興兒聽見這個聲音兒，早已沒了主意了，只得乍著膽子進來。鳳姐兒一見，便說：「好小子啊！你和你爺辦的好事啊！你只實說罷！」興兒一聞此言，又看見鳳姐氣色，早唬軟了，不覺跪下，只是磕頭。鳳姐兒道：「論起這事來，我也聽見說不與你相干。但只你不早來回我知道，這就是你的不是了。你要實說了，我還饒你；再有一字虛言，你先摸摸你腔子上幾個腦袋瓜子！」興兒戰戰兢兢的朝上磕頭道：「奶奶問的是什麼事，奴才同爺辦壞了？」鳳姐聽了，一腔火都發作起來，喝命：「打嘴巴！」旺兒過來才要打時，鳳姐兒罵道：「什麼糊塗忘八崽子！叫他自己打，用你打嗎！一會子你再各人打你那嘴巴子還不遲呢。」那興兒真個自己左右開弓打了自己十幾個嘴巴。鳳姐兒喝聲「站住」，問道：「你二爺外頭娶了什麼新奶奶舊奶奶的事，你大概不知道啊。」興兒見說出這件事來，越發著了慌，連忙把帽子抓下來在磚地上咕咚咕咚碰的頭山響，口裡說道：「只求奶奶超生，奴才再不敢撒一個字兒的謊。」鳳姐道：「快說！」興兒直蹶蹶的跪起來回道：「這事頭裡奴才也不知道。就是這一天，東府里大老爺送了殯，俞祿往珍大爺廟裡去領銀子。二爺同著蓉哥兒到了東府里，道兒上爺兒兩個說起珍大奶奶那邊的二位姨奶奶來。二爺誇他好，蓉哥兒哄著二爺，說把二姨奶奶說給二爺。」鳳姐聽到這裡，使勁啐道：「呸，沒臉的忘八蛋！他是你那一門子的姨奶奶！」興兒忙又磕頭說：「奴才該死！」往上啾著，不敢言語。鳳姐兒道：「完了嗎？怎麼不說了？」興兒方纔又回道：「奶奶恕奴才，奴才才敢回。」鳳姐啐道：「放你媽的屁，這還什麼恕不恕了。你好生給我往下說，好多著呢。」 興兒又回道：「二爺聽見這個話就喜歡了。後來奴才也不知道怎麼就弄真了。」鳳姐微微冷笑道：「這個自然麽，你可那裡知道呢！你知道的只怕都煩了呢。是了，說底下的罷！」興兒回道：「後來就是蓉哥兒給二爺找了房子。」鳳姐忙問道：「如今房子在那裡？」興兒道：「就在府後頭。」鳳姐兒道：「哦。」回頭瞅著平兒道：「咱們都是死人哪。你聽聽！」平兒也不敢作聲。興兒又回道：「珍大爺那邊給了張家不知多少銀子，那張家就不問了。」鳳姐道：「這裡頭怎麼又扯拉上什麼張家李家咧呢？」興兒回道：「奶奶不知道，這二奶奶……」剛說到這裡，又自己打了個嘴巴，想了想，說道： 「那珍大奶奶的妹子……」鳳姐兒接著道：「怎麼樣？快說呀。」興兒道：「那珍大奶奶的妹子原來從小兒有人家的，姓張，叫什麼張華，如今窮的待好討飯。珍大爺許了他銀子，他就退了親了。」鳳姐兒聽到這裡，點了點頭兒，回頭便望平兒說道：「你都聽見了？小忘八崽子，頭裡他還說他不知道呢！」興兒又回道： 「後來二爺才叫人裱糊了房子，娶過來了。」鳳姐道：「打那裡娶過來的？」興兒回道：「就在他老娘家抬過來的。」鳳姐又問：「沒人送親麽？」興兒道：「就是蓉哥兒。還有幾個丫頭老婆子們，沒別人。」鳳姐道：「你大奶奶沒來嗎？」興兒道：「過了兩天，大奶奶才拿了些東西來瞧的。」鳳姐兒回頭向平兒道：「怪道那兩天二爺稱贊大奶奶不離嘴呢。」掉過臉來又問興兒，「誰伏侍呢？自然是你了。」興兒趕著碰頭不言語。鳳姐又問：「前頭那些日子說給那府里辦事，想來辦的就是這個了。」興兒回道：「也有辦事的時候，也有往新房子里去的時候。」

鳳姐聽了這一篇言詞，只氣得癡呆了半天，面如金紙，兩隻吊稍丹鳳眼越發直豎起來了，渾身亂戰。半晌，連話也說不上來，只是發怔。猛低頭，見興兒還在地下跪着，便說道：「你這個猴兒崽子就該打死。這有什麼瞞著我的？你想著瞞了我，就在你那糊塗爺跟前討了好兒了，你新奶奶好疼你。」興兒道：「未能早回奶奶，是奴才該死！」便叩頭有聲。

鳳姐又問道：「誰和他住著呢。」興兒道：「先是和他娘和妹子在一處。就在十幾天前，他妹子自己抹了脖子。他娘得病，昨兒也死了。」鳳姐道：「這又為什麼？」興兒隨將柳湘蓮的事說了一遍。鳳姐道：「這個人還算造化高，省了當那出名兒的忘八。」因又問道：「沒了別的事了麽？」興兒道：「別的事奴才不知道。奴才剛纔說的字字是實話，一字虛假，奶奶問出來只管打死奴才，奴才也無怨的。」鳳姐低了一回頭，便又指著興兒說道：「我不看你剛纔還有點怕懼兒，不敢撒謊，我把你的腿不給你砸折了。」說著喝聲「出去！」興兒瞌了個頭，才爬起來，退到外間門口，不敢就走。鳳姐道：「過來，我還有話呢。」興兒趕忙垂手敬聽。鳳姐道：「你忙什麼，新奶奶等著賞你什麼呢？」興兒也不敢抬頭。鳳姐道：「我什麼時候叫你，你什麼時候到。遲一步兒，你試試！出去罷。」興兒忙答應幾個「是」，退出門來。鳳姐又叫道：「興兒！」興兒趕忙站住。鳳姐道：「快出去告訴你二爺去，是不是啊？」興兒回道：「奴才不敢。」鳳姐道：「你出去提一個字兒，隄防你的皮！」興兒連忙答應著才出去了。鳳姐又叫：「旺兒呢？」旺兒連忙答應著過來。鳳姐把眼直瞪瞪的瞅了兩三句話的工夫，才說道：「好旺兒，很好，去罷！外頭有人提一個字兒，全在你身上。」旺兒答應著也出去了。

且說鳳姐見興兒出去，回頭向平兒說：「方才興兒說的話，你都聽見了沒有？天下那有這樣沒臉的男人！吃着碗裏，看着鍋里，見一個，愛一個，真成了餵不飽的狗，實在是個棄舊迎新的壞貨。只可惜這五六品的頂帶給他！他別想着俗語說的『家花那有野花香』的話，他要信了這個話，可就大錯了。多早晚在外面鬧一個沒臉、親戚朋友見不得的事出來，他才罷手呢！」平兒一旁勸道：「奶奶身子才好了，也不可過於氣惱。看二爺自從鮑二的女人那一件事之後，倒收了心，好了呢，如今為什麼又干起這樣事來？這都是珍大爺他的不是。」鳳姐說：「珍大爺固有不是，也總因咱們那位下作不堪的爺他眼饞，人家才引誘他的。俗語說『牛兒不吃水，也強按頭麼？』珍大爺幹這樣的事，珍大奶奶也該攔着不依才是。珍大奶奶也不想一想，把一個妹子要許幾家子弟才好，先許了姓張的，今又嫁了姓賈的；天下的男人都死絕了，都嫁到賈家來！難道賈家的衣食這樣好不成？那妹子本來也不是他親的，而且聽見說原是個混賬爛桃。難道珍大奶奶現做着命婦，家中有這樣一個打嘴現世的妹子，也不知道羞臊，躲避着些，反倒大面上揚鈴打鼓的，在這門裏丟醜，也不怕笑話？珍大爺也是做官的人，別的律例不知道也罷了，連個服中娶親，停妻再娶，使不得的規矩，他也不知道不成？你替他細想一想，他幹的這件事，是疼兄弟，還是害兄弟呢？」平兒說：「珍大爺只顧眼前，叫兄弟喜歡，也不管日後的輕重干係了。」鳳姐兒冷笑道：「這是什麼『叫兄弟喜歡』，這是給他毒藥吃！若論親叔伯兄弟中，他年紀又最大，又居長，不知教導兄弟學好，反引誘兄弟學不長進，擔罪名兒，日後鬧出事來，他在一邊缸沿兒上站着看熱鬧，真真我要罵也罵不出口來。他在那邊府里的醜事壞名聲，已經叫人聽不上了，必定也叫兄弟學他一樣，才好顯不出他的丑來。這是什麼作哥哥的道理？倒不如撒泡尿浸死了，替大老爺死了也罷，活着作什麼。」

平兒看鳳姐越說越氣，便跪在地下，再三苦勸安慰一會子，鳳姐才略消了些氣惱。喝了口茶，喘息一回，便要了拐枕，歪在牀上，閉眼養神。平兒只得悄悄的退出去了。鳳姐將前事從頭至尾細細的盤算多時，才得了主意，也不吿訴平兒，卻作出個嘻笑自若、毫無惱恨妒嫉的樣子來。心下早已算定，只待賈璉起程去平安州，再作道理。要知端的，且聽下回分解。
